\documentclass[12pt,a4paper]{article}

\input{/home/tabaka/Documents/GitHub/Defs/defs.tex}

\setcounter{zadaniaRozdzial}{3}
\setcounter{section}{0}

\begin{document}
	
	\begin{center}
		\Large Wymagania edukacyjne z matematyki w klasie 1 Szkoły Branżowej na poszczególne oceny.
	\end{center}
	
	Wymagania edukacyjne na poszczególne oceny są stopniowane.
	Uzyskanie danej oceny wymaga spełnienia wymagań przewidzianych
	dla wszystkich ocen niższych, z uwzględnieniem wymagań
	określonych dla danej oceny.
	
	\section{Liczby rzeczywiste}
	
	\begin{longtable}{|p{16cm}|}
		\hline
		
		\textbf{Wymagania edukacyjne na ocenę dopuszczającą} \\ \hline
		
		\textbullet\ używa powszechnie przyjętych oznaczeń zbiorów liczbowych: $\mathbb{N}$ – liczb naturalnych, $\mathbb{Z}$ – liczb całkowitych, $\mathbb{Q}$ – liczb wymiernych i $\mathbb{R}$ – liczb rzeczywistych; \\
		
		\textbullet\ zna cechy podzielności liczb przez $2$, $3$, $5$ i $9$; \\
		
		\textbullet\ stosuje cechy podzielności liczb przez $2$ i $5$; \\
		
		\textbullet\ zna pojęcie dzielnika liczby naturalnej; \\
		
		\textbullet\ podaje dzielniki liczb naturalnych w prostych przypadkach; \\
		
		\textbullet\ wykonuje dzielenie liczby naturalnej z resztą; \\
		
		\textbullet\ rozpoznaje wśród podanych liczb liczby naturalne, całkowite i wymierne; \\
		
		\textbullet\ zna zasady kolejności wykonywania działań; \\
		
		\textbullet\ posługuje się kalkulatorem podczas wykonywania obliczeń; \\
		
		\textbullet\ wykonuje proste działania na liczbach wymiernych; \\
		
		\textbullet\ wyznacza rozwinięcia dziesiętne ułamków zwykłych; \\
		
		\textbullet\ zamienia ułamki dziesiętne o skończonym rozwinięciu na ułamki zwykłe; \\
		
		\textbullet\ zna regułę zaokrąglania liczb. \\ \hline
		
		\textbf{Wymagania edukacyjne na ocenę dostateczną} \\ \hline
		
		\textbullet\ stosuje cechy podzielności liczb przez $3$ i $9$; \\
		
		\textbullet\ wypisuje dzielniki liczby naturalnej; \\
		
		\textbullet\ stosuje działania na liczbach naturalnych w sytuacjach praktycznych; \\
		
		\textbullet\ stosuje zasady kolejności wykonywania działań w prostych przypadkach; \\
		
		\textbullet\ stosuje regułę zaokrąglania liczb w prostych przypadkach; \\
		
		\textbullet\ wyznacza wskazaną cyfrę po przecinku w rozwinięciu dziesiętnym okresowym liczby; \\
		
		\textbullet\ wykonuje działania na liczbach wymiernych; \\
		
		\textbullet\ szacuje wyniki działań; \\
		
		\textbullet\ zna pojęcie błędu przybliżenia; \\
		
		\textbullet\ zna pojęcie wartości bezwzględnej. \\ \hline
		
		\textbf{Wymagania edukacyjne na ocenę dobrą} \\ \hline
		
		\textbullet\ stosuje działania na liczbach całkowitych w sytuacjach praktycznych; \\
		
		\textbullet\ oblicza wartości wyrażeń arytmetycznych, stosując zasady kolejności wykonywania działań; \\
		
		\textbullet\ zna pojęcia dziennego zapotrzebowania energetycznego, wartości energetycznej produktu i wartości kalorycznej porcji; \\
		
		\textbullet\ zaokrągla liczbę z podaną dokładnością; \\
		
		\textbullet\ oblicza wartości wyrażeń zawierających wartość bezwzględną; \\
		
		\textbullet\ posługuje się rozwinięciem dziesiętnym liczby w rozliczeniach finansowych; \\
		
		\textbullet\ oblicza wartość bezwzględną liczby wymiernej; \\
		
		\textbullet\ oblicza błąd przybliżenia danej liczby oraz określa, czy przybliżenie jest z nadmiarem, czy z niedomiarem. \\ \hline
		
		\textbf{Wymagania edukacyjne na ocenę bardzo dobrą} \\ \hline
		
		\textbullet\ stosuje działania na liczbach wymiernych w sytuacjach praktycznych; \\
		
		\textbullet\ posługuje się pojęciami dziennego zapotrzebowania energetycznego, wartości energetycznej produktu i wartości kalorycznej porcji; \\
		
		\textbullet\ oblicza wartość energetyczną posiłków; \\
		
		\textbullet\ oblicza błąd bezwzględny przybliżenia danej liczby; \\
		
		\textbullet\ stosuje przybliżenia w sytuacjach praktycznych. \\ \hline
		
		\textbf{Wymagania edukacyjne na ocenę celującą} \\ \hline
		
		\textbullet\ rozwiązuje zadania o podwyższonym stopniu trudności. \\ \hline
		
	\end{longtable}
	
	\section{Potęgi i pierwiastki}
	
	\begin{longtable}{|p{16cm}|}
		\hline
		
		\textbf{Wymagania edukacyjne na ocenę dopuszczającą} \\ \hline
		
		\textbullet\ zna pojęcie potęgi; \\
		
		\textbullet\ oblicza wartości potęg o wykładnikach naturalnych; \\
		
		\textbullet\ zna pojęcia pierwiastka kwadratowego i pierwiastka sześciennego; \\
		
		\textbullet\ oblicza wartość pierwiastka kwadratowego i sześciennego z liczby nieujemnej; \\
		
		\textbullet\ zna prawa działań na potęgach i pierwiastkach; \\
		
		\textbullet\ zna pojęcie liczby niewymiernej. \\ \hline
		
		\textbf{Wymagania edukacyjne na ocenę dostateczną} \\ \hline
		
		\textbullet\ oblicza wartości potęg o wykładnikach całkowitych ujemnych i podstawach wymiernych w prostych przypadkach; \\
		
		\textbullet\ oblicza wartość pierwiastka sześciennego z liczby ujemnej; \\
		
		\textbullet\ stosuje potęgi o podstawie $10$ przy zamianie jednostek długości i pola powierzchni; \\
		
		\textbullet\ stosuje prawa działań na potęgach i pierwiastkach do upraszczania wyrażeń lub obliczania ich wartości w prostych przypadkach; \\
		
		\textbullet\ szacuje wartości pierwiastków. \\ \hline
		
		\textbf{Wymagania edukacyjne na ocenę dobrą} \\ \hline
		
		\textbullet\ oblicza wartości potęg o wykładnikach całkowitych ujemnych i podstawach wymiernych; \\
		
		\textbullet\ stosuje prawa działań na potęgach i pierwiastkach do upraszczania wyrażeń lub obliczania ich wartości; \\
		
		\textbullet\ wyłącza czynnik przed znak pierwiastka; \\
		
		\textbullet\ oblicza przybliżone wartości liczb niewymiernych postaci $a\sqrt{b}$. \\ \hline
		
		\textbf{Wymagania edukacyjne na ocenę bardzo dobrą} \\ \hline
		
		\textbullet\ stosuje działania na potęgach i pierwiastkach w sytuacjach praktycznych; \\
		
		\textbullet\ oblicza wartość pierwiastka stopnia wyższego niż trzeci. \\ \hline
		
		\textbf{Wymagania edukacyjne na ocenę celującą} \\ \hline
		
		\textbullet\ rozwiązuje zadania o podwyższonym stopniu trudności. \\ \hline
		
	\end{longtable}
	
	\section{Procenty i ich zastosowanie}
	
	\begin{longtable}{|p{16cm}|}
		\hline
		
		\textbf{Wymagania edukacyjne na ocenę dopuszczającą} \\ \hline
		
		\textbullet\ zna pojęcia procentu i punktu procentowego; \\
		
		\textbullet\ zamienia procenty na ułamki i odwrotnie; \\
		
		\textbullet\ oblicza procent danej liczby w prostych przypadkach; \\
		
		\textbullet\ zna pojęcia kwoty/ceny netto, kwoty/ceny brutto oraz podatku VAT; \\
		
		\textbullet\ odczytuje informacje z faktury VAT. \\ \hline
		
		\textbf{Wymagania edukacyjne na ocenę dostateczną} \\ \hline
		
		\textbullet\ oblicza procent danej liczby; \\
		
		\textbullet\ zwiększa i zmniejsza liczbę o dany procent; \\
		
		\textbullet\ oblicza, jakim procentem jednej liczby jest druga liczba; \\
		
		\textbullet\ wyznacza liczbę, gdy dany jest jej procent, w prostych przypadkach; \\
		
		\textbullet\ oblicza kwotę/cenę brutto lub podatek VAT, gdy dana jest kwota/cena netto. \\ \hline
		
		\textbf{Wymagania edukacyjne na ocenę dobrą} \\ \hline
		
		\textbullet\ wyznacza liczbę, gdy dany jest jej procent; \\
		
		\textbullet\ oblicza jedną z wielkości: podatek VAT, kwotę/cenę netto lub kwotę/cenę brutto, gdy dane są dwie pozostałe; \\
		
		\textbullet\ zna pojęcia procentu składanego i kapitalizacji odsetek; \\
		
		\textbullet\ stosuje obliczenia procentowe w prostych zadaniach praktycznych, np. do obliczania kosztów kredytów lub zysków z lokat. \\ \hline
		
		\textbf{Wymagania edukacyjne na ocenę bardzo dobrą} \\ \hline
		
		\textbullet\ stosuje obliczenia procentowe i własności potęgowania do obliczania zysków z lokat oprocentowanych procentem składanym, z uwzględnieniem kapitalizacji odsetek w okresach krótszych niż rok; \\
		
		\textbullet\ analizuje różne oferty kredytowe i wybiera korzystniejszą z nich. \\ \hline
		
		\textbf{Wymagania edukacyjne na ocenę celującą} \\ \hline
		
		\textbullet\ rozwiązuje zadania o podwyższonym stopniu trudności. \\ \hline
		
	\end{longtable}
	
	\section{Równania i nierówności}
	
	\begin{longtable}{|p{16cm}|}
		\hline
		
		\textbf{Wymagania edukacyjne na ocenę dopuszczającą} \\ \hline
		
		\textbullet\ zna pojęcia równania równoważnego, sprzecznego i tożsamościowego; \\
		
		\textbullet\ sprawdza, czy dana liczba jest rozwiązaniem równania; \\
		
		\textbullet\ zna metodę równań równoważnych; \\
		
		\textbullet\ odczytuje z osi liczbowej współrzędną danego punktu oraz zaznacza na osi liczbowej punkt o danej współrzędnej; \\
		
		\textbullet\ zaznacza na osi liczbowej przedziały opisane symbolicznie lub za pomocą nierówności; \\
		
		\textbullet\ rozróżnia przedziały otwarte, domknięte, lewostronnie domknięte, prawostronnie domknięte i nieograniczone; \\
		
		\textbullet\ zna pojęcia nierówności ostrej i nieostrej. \\ \hline
		
		\textbf{Wymagania edukacyjne na ocenę dostateczną} \\ \hline
		
		\textbullet\ rozpoznaje równania sprzeczne i tożsamościowe oraz podaje ich zbiory rozwiązań; \\
		
		\textbullet\ stosuje przekształcenia równoważne do rozwiązywania równań; \\
		
		\textbullet\ odczytuje i zapisuje symbolicznie lub za pomocą nierówności przedział zaznaczony na osi liczbowej; \\
		
		\textbullet\ sprawdza, czy dana liczba jest rozwiązaniem nierówności; \\
		
		\textbullet\ sprawdza, czy nierówności są równoważne; \\
		
		\textbullet\ stosuje przekształcenia równoważne do rozwiązywania prostych nierówności pierwszego stopnia z jedną niewiadomą; \\
		
		\textbullet\ zapisuje zbiór rozwiązań nierówności w postaci przedziału. \\ \hline
		
		\textbf{Wymagania edukacyjne na ocenę dobrą} \\ \hline
		
		\textbullet\ stosuje równania pierwszego stopnia z jedną niewiadomą do rozwiązywania prostych zadań tekstowych; \\
		
		\textbullet\ stosuje przekształcenia równoważne do rozwiązywania nierówności pierwszego stopnia z jedną niewiadomą; \\
		
		\textbullet\ rozpoznaje nierówności sprzeczne i tożsamościowe oraz podaje ich zbiory rozwiązań. \\ \hline
		
		\textbf{Wymagania edukacyjne na ocenę bardzo dobrą} \\ \hline
		
		\textbullet\ stosuje równania i nierówności pierwszego stopnia z jedną niewiadomą do rozwiązywania zadań osadzonych w sytuacjach praktycznych. \\ \hline
		
		\textbf{Wymagania edukacyjne na ocenę celującą} \\ \hline
		
		\textbullet\ rozwiązuje zadania o podwyższonym stopniu trudności. \\ \hline
		
	\end{longtable}
	
	\section{Układy równań}
	
	\begin{longtable}{|p{16cm}|}
		\hline
		
		\textbf{Wymagania edukacyjne na ocenę dopuszczającą} \\ \hline
		
		\textbullet\ zna pojęcie układu równań liniowych z dwiema niewiadomymi; \\
		
		\textbullet\ rozróżnia metody rozwiązywania układów równań; \\
		
		\textbullet\ sprawdza, czy dana para liczb jest rozwiązaniem układu równań liniowych z dwiema niewiadomymi; \\
		
		\textbullet\ zna pojęcia układu równań oznaczonego, nieoznaczonego i sprzecznego. \\ \hline
		
		\textbf{Wymagania edukacyjne na ocenę dostateczną} \\ \hline
		
		\textbullet\ określa, czy dany układ równań jest oznaczony, nieoznaczony czy sprzeczny; \\
		
		\textbullet\ rozwiązuje układy równań metodą podstawiania lub metodą przeciwnych współczynników w prostych przypadkach. \\ \hline
		
		\textbf{Wymagania edukacyjne na ocenę dobrą} \\ \hline
		
		\textbullet\ rozwiązuje układy równań metodą podstawiania lub metodą przeciwnych współczynników; \\
		
		\textbullet\ układa i rozwiązuje układy równań do prostych zadań tekstowych. \\ \hline
		
		\textbf{Wymagania edukacyjne na ocenę bardzo dobrą} \\ \hline
		
		\textbullet\ układa i rozwiązuje układy równań do zadań tekstowych osadzonych w sytuacjach praktycznych. \\ \hline
		
		\textbf{Wymagania edukacyjne na ocenę celującą} \\ \hline
		
		\textbullet\ rozwiązuje zadania o podwyższonym stopniu trudności. \\ \hline
		
	\end{longtable}
	
	\section{Funkcje}
	
	\begin{longtable}{|p{16cm}|}
		\hline
		
		\textbf{Wymagania edukacyjne na ocenę dopuszczającą} \\ \hline
		
		\textbullet\ zna pojęcia funkcji, argumentu, dziedziny, wartości funkcji i miejsca zerowego; \\
		
		\textbullet\ rozpoznaje wśród podanych przyporządkowań te, które są funkcjami; \\
		
		\textbullet\ zna różne sposoby przedstawiania funkcji: opis słowny, graf, tabelę, wzór i wykres; \\
		
		\textbullet\ zaznacza punkty w układzie współrzędnych na płaszczyźnie; \\
		
		\textbullet\ odczytuje współrzędne danych punktów; \\
		
		\textbullet\ zna pojęcia funkcji stałej, rosnącej i malejącej. \\ \hline
		
		\textbf{Wymagania edukacyjne na ocenę dostateczną} \\ \hline
		
		\textbullet\ rysuje wykres funkcji liczbowej określonej prostym wzorem; \\
		
		\textbullet\ oblicza wartość funkcji dla danego argumentu na podstawie jej wzoru; \\
		
		\textbullet\ odczytuje z wykresu dziedzinę i zbiór wartości funkcji, jej miejsca zerowe, argumenty, dla których przyjmuje ona daną wartość, oraz argumenty, dla których przyjmuje wartość największą lub najmniejszą w danym przedziale; \\
		
		\textbullet\ rozpoznaje wśród podanych wykresów wykresy funkcji monotonicznych: rosnących, malejących i stałych; \\
		
		\textbullet\ na podstawie wykresu funkcji $y=f(x)$ rysuje wykresy funkcji $y=f(x)+q$ dla $q>0$ oraz $q<0$. \\ \hline
		
		\textbf{Wymagania edukacyjne na ocenę dobrą} \\ \hline
		
		\textbullet\ rysuje wykres funkcji liczbowej określonej wzorem; \\
		
		\textbullet\ odczytuje z wykresu przedziały, w których funkcja przyjmuje wartości dodatnie lub ujemne oraz wartości mniejsze lub większe od podanej liczby; \\
		
		\textbullet\ określa na podstawie wykresu przedziały monotoniczności funkcji; \\
		
		\textbullet\ stosuje własności funkcji monotonicznej do porównywania jej wartości dla danych argumentów; \\
		
		\textbullet\ na podstawie wykresu funkcji $y=f(x)$ rysuje wykresy funkcji $y=f(x-p)$ dla $p>0$ oraz $p<0$; \\
		
		\textbullet\ rozpoznaje zależności funkcyjne w zadaniach osadzonych w kontekście praktycznym. \\ \hline
		
		\textbf{Wymagania edukacyjne na ocenę bardzo dobrą} \\ \hline
		
		\textbullet\ przedstawia za pomocą wykresu funkcję liczbową określoną różnymi wzorami na różnych przedziałach; \\
		
		\textbullet\ rozwiązuje zadania tekstowe z wykorzystaniem funkcji przedstawionych w różnych postaciach, np. za pomocą wykresu lub wzoru; \\
		
		\textbullet\ przedstawia zależności funkcyjne opisane w zadaniach tekstowych za pomocą wzoru lub wykresu; \\
		
		\textbullet\ stosuje monotoniczność funkcji w zadaniach osadzonych w sytuacjach praktycznych. \\ \hline
		
		\textbf{Wymagania edukacyjne na ocenę celującą} \\ \hline
		
		\textbullet\ rozwiązuje zadania o podwyższonym stopniu trudności. \\ \hline
		
	\end{longtable}
	
	\section{Funkcja liniowa}
	
	\begin{longtable}{|p{16cm}|}
		\hline
		
		\textbf{Wymagania edukacyjne na ocenę dopuszczającą} \\ \hline
		
		\textbullet\ zna wzór funkcji liniowej oraz pojęcie współczynnika kierunkowego; \\
		
		\textbullet\ rozpoznaje wśród podanych wzorów funkcji wzór funkcji liniowej; \\
		
		\textbullet\ sprawdza, czy dany punkt należy do wykresu funkcji liniowej określonej podanym wzorem; \\
		
		\textbullet\ zna warunek równoległości wykresów funkcji liniowych; \\
		
		\textbullet\ odczytuje z wykresu funkcji współrzędne punktów przecięcia prostej z osiami układu współrzędnych; \\
		
		\textbullet\ zna pojęcia kosztów produkcji, przychodu, zysku i straty. \\ \hline
		
		\textbf{Wymagania edukacyjne na ocenę dostateczną} \\ \hline
		
		\textbullet\ rysuje wykres funkcji liniowej na podstawie jej wzoru; \\
		
		\textbullet\ wyznacza współrzędne punktów przecięcia prostej danej równaniem kierunkowym z osiami układu współrzędnych; \\
		
		\textbullet\ interpretuje współczynniki występujące we wzorze funkcji liniowej, wskazując funkcje, których wykresy są prostymi równoległymi; \\
		
		\textbullet\ określa monotoniczność funkcji liniowej danej wzorem; \\
		
		\textbullet\ posługuje się pojęciami kosztów produkcji, przychodu, zysku i straty. \\ \hline
		
		\textbf{Wymagania edukacyjne na ocenę dobrą} \\ \hline
		
		\textbullet\ rysuje prostą, wykorzystując interpretację geometryczną współczynnika kierunkowego; \\
		
		\textbullet\ oblicza współczynnik kierunkowy prostej na podstawie jej wykresu lub współrzędnych dwóch punktów należących do tej prostej; \\
		
		\textbullet\ rozwiązuje układy równań metodą graficzną. \\ \hline
		
		\textbf{Wymagania edukacyjne na ocenę bardzo dobrą} \\ \hline
		
		\textbullet\ wyznacza wzór funkcji liniowej na podstawie informacji o tej funkcji lub jej wykresie; \\
		
		\textbullet\ wykorzystuje zależność między liczbą rozwiązań układu równań a wzajemnym położeniem wykresów odpowiadających mu prostych; \\
		
		\textbullet\ analizuje wyniki finansowe firmy w przypadkach, gdy jej przychód i koszty są opisane za pomocą wykresów lub wzorów. \\ \hline
		
		\textbf{Wymagania edukacyjne na ocenę celującą} \\ \hline
		
		\textbullet\ rozwiązuje zadania o podwyższonym stopniu trudności. \\ \hline
		
	\end{longtable}
	
	\section{Statystyka}
	
	\begin{longtable}{|p{16cm}|}
		\hline
		
		\textbf{Wymagania edukacyjne na ocenę dopuszczającą} \\ \hline
		
		\textbullet\ zna pojęcia średniej arytmetycznej, średniej ważonej, mediany i dominanty; \\
		
		\textbullet\ oblicza średnią arytmetyczną podanych liczb; \\
		
		\textbullet\ wyznacza dominantę. \\ \hline
		
		\textbf{Wymagania edukacyjne na ocenę dostateczną} \\ \hline
		
		\textbullet\ oblicza średnią arytmetyczną danych przedstawionych w tabeli lub na diagramie; \\
		
		\textbullet\ oblicza średnią ważoną liczb o podanych wagach; \\
		
		\textbullet\ wyznacza medianę zestawu danych. \\ \hline
		
		\textbf{Wymagania edukacyjne na ocenę dobrą} \\ \hline
		
		\textbullet\ wyznacza medianę danych przedstawionych w tabeli lub na diagramie; \\
		
		\textbullet\ wykorzystuje średnią arytmetyczną, średnią ważoną, medianę i dominantę do rozwiązywania prostych zadań tekstowych. \\ \hline
		
		\textbf{Wymagania edukacyjne na ocenę bardzo dobrą} \\ \hline
		
		\textbullet\ wykorzystuje miary statystyczne do rozwiązywania zadań osadzonych w kontekście praktycznym. \\ \hline
		
		\textbf{Wymagania edukacyjne na ocenę celującą} \\ \hline
		
		\textbullet\ rozwiązuje zadania o podwyższonym stopniu trudności. \\ \hline
		
	\end{longtable}
	
\end{document}